\documentclass[12pt,a4paper]{article}
\usepackage[utf8]{inputenc}
\usepackage[T2A]{fontenc}
\usepackage[russian]{babel}
\usepackage{amsmath,amssymb}
\usepackage{array}
\usepackage{longtable}
\usepackage{booktabs}
\usepackage{geometry}
\usepackage{hyperref}
\usepackage{graphicx}
\usepackage{xurl}
\geometry{margin=2.3cm}

\title{Единый спектральный источник для \texttt{TOE.pdf} и \texttt{UGSM}:\\
сводная статья по мосту, численному аудиту и углублённой проверке гипотезы общей тени}
\author{}
\date{}

\begin{document}

\maketitle

\begin{abstract}
В статье исследуется гипотеза, что теоретические корпуса \path{TOE.pdf} и \path{ugsm_dynamics_audit-3.pdf} можно интерпретировать как две разные тени одной более общей спектральной сущности. Работа объединяет оба выполненных цикла исследования в один документ: концептуальный мост между корреляционным объектом TOE и спектральным действием UGSM, первичный численный аудит на \(S^3\times S^1\), а затем углублённую проверку устойчивости по окну фита, радиусам и семейству спектральных kernels. Показано, что минимальный тепловой объект на \(S^3\times S^1\) воспроизводит ключевые геометрические масштабы \(4\pi^3\), \(\pi^2\) и естественный \(\pi\)-proxy, а сама картина сохраняется при нескольких независимых деформациях вычислительной схемы. Итог статьи --- не теорема о полном тождестве двух теорий, а существенно усиленная и численно поддержанная гипотеза единого спектрального источника.
\end{abstract}

\section{Введение}

Сравнение \texttt{TOE.pdf} и \texttt{ugsm\_dynamics\_audit-3.pdf} естественно начинается с различий в масштабе. TOE претендует на широкую архитектурную унификацию, тогда как UGSM организована как более локальный, но процедурно дисциплинированный спектральный EFT-аудит. Однако за различием масштаба может скрываться более глубокий вопрос: не являются ли обе конструкции двумя разными функциональными чтениями одного и того же более общего спектрального уровня.

Именно этот вопрос и стал предметом настоящего исследования. Важно подчеркнуть, что задача статьи не в том, чтобы риторически сблизить два корпуса, а в том, чтобы:
\begin{itemize}
    \item дать минимальный математический мост между ними;
    \item проверить этот мост на простейшем контролируемом компактном фоне;
    \item затем проверить, остаётся ли найденная картина устойчивой при дополнительных численных стресс-тестах.
\end{itemize}

\section{Исследовательский вопрос и гипотеза}

Основной исследовательский вопрос формулируется так:
\begin{quote}
существует ли единый базовый операторный источник, из которого корреляционный объект TOE и спектральный функционал UGSM получаются как две разные функциональные репрезентации?
\end{quote}

Рабочая гипотеза статьи имеет вид
\begin{equation}
\hat C_\sigma \sim e^{-\sigma D_M^2/2},
\end{equation}
где \(\hat C_\sigma\) --- корреляционный объект TOE, а \(D_M^2\) --- оператор Лаплас/Дирак-типа на компактном секторе UGSM.

Если функция \(f\) в UGSM допускает представление Лапласа
\begin{equation}
f(x)=\int_0^{\infty}\mu(t)e^{-tx}\,dt,
\end{equation}
то спектральное действие UGSM можно записать как
\begin{equation}
S_{spec}(\Lambda)=\mathrm{Tr}\,f(D_M^2/\Lambda^2)
=\int_0^{\infty}\mu(t)\,\mathrm{Tr}\,e^{-tD_M^2/\Lambda^2}\,dt.
\end{equation}

После замены \(\sigma=2t/\Lambda^2\) получается минимальный мост:
\begin{equation}
S_{spec}(\Lambda)=\int_0^{\infty}\mu(t)\,\mathrm{Tr}\,\hat C_{2t/\Lambda^2}\,dt.
\end{equation}

Следовательно, UGSM может быть прочитана как спектрально-моментная редукция того же источника, который в TOE считается фундаментальным. Именно в этом смысле и используется выражение ``общая тень''.

\section{Материалы и метод}

Для численной проверки был выбран минимальный компактный сектор
\[
M=S^3(R_3)\times S^1(R_1),
\]
который естественно согласуется с UGSM-постановкой. На первом этапе использовался скалярный лапласиан с собственными значениями и кратностями
\begin{align}
\lambda_{n,m} &= \frac{n(n+2)}{R_3^2}+\frac{m^2}{R_1^2},\\
 g_n &= (n+1)^2.
\end{align}

Тепловой след определялся формулой
\begin{equation}
K(t)=\sum_{n,m} g_n e^{-t\lambda_{n,m}},
\end{equation}
а анализировалась стандартная нормированная величина
\begin{equation}
\widetilde K(t)=(4\pi t)^2 K(t).
\end{equation}

Далее она аппроксимировалась малой асимптотикой
\begin{equation}
\widetilde K(t)\approx a_0+a_2 t+a_4 t^2.
\end{equation}

Исследование состояло из двух циклов:
\begin{enumerate}
    \item \textbf{Первичный цикл:} базовый fit на \(R_3=R_1=1\), проверка коэффициентов \(a_0\), \(a_2\), \(a_2/(4\pi)\) и интерпретация \(\pi\)-сектора.
    \item \textbf{Углублённый цикл:} проверка устойчивости по окну фита, радиусам \((R_3,R_1)\) и семейству спектральных kernels.
    \item \textbf{Dirac-type цикл:} перенос схемы в спинорный сектор для \(D^2\) на \(S^3\times S^1\), проверка спинорных коэффициентов и устойчивости моста в более физически содержательной постановке.
    \item \textbf{Spin-holonomy цикл:} проверка twisted spin-структур и Wilson/holonomy-сдвига \(\alpha\) на \(S^1\), чтобы отделить локальные Seeley--DeWitt коэффициенты от глобального \(\pi\)-сектора.
    \item \textbf{Gauge-holonomy цикл:} первый шаг к более полным Wilson-loop деформациям: введение абелевого gauge-holonomy параметра \(\beta\) и внутреннего doublet-сектора с зарядами \(q=\pm1\).
    \item \textbf{Enriched-operator цикл:} шестой цикл вводит минимальный richer Dirac-type оператор с внутренним charge/mass multiplet, чтобы проверить, сохраняется ли мост при появлении нескольких внутренних каналов и субведущей структурной нагрузки.
    \item \textbf{Sector-map цикл:} седьмой цикл строит карту вкладов по параметрам \(\mu\) и \(\beta\), чтобы отделить mass-driven subleading-структуру от чисто holonomy-driven фазового сектора.
    \item \textbf{Sector-attribution цикл:} восьмой цикл добавляет явную sensitivity-декомпозицию и словарь соответствий между операторными компонентами и возможными блоками TOE/UGSM.
\end{enumerate}

Первый цикл реализован в файлах \path{spectral_bridge_experiment.py} и \path{spectral_bridge_results.json}. Второй цикл реализован в файлах \path{spectral_unity_deep_checks.py} и \path{spectral_unity_deep_results.json}. Третий цикл реализован в файлах \path{dirac_unity_audit.py} и \path{dirac_unity_results.json}. Четвёртый цикл реализован в файлах \path{dirac_spin_holonomy_audit.py} и \path{dirac_spin_holonomy_results.json}. Пятый цикл реализован в файлах \path{gauge_holonomy_audit.py} и \path{gauge_holonomy_results.json}. Шестой цикл реализован в файлах \path{enriched_dirac_operator_audit.py} и \path{enriched_dirac_operator_results.json}. Седьмой цикл реализован в файлах \path{enriched_sector_map_audit.py} и \path{enriched_sector_map_results.json}. Восьмой цикл реализован в файлах \path{sector_attribution_audit.py} и \path{sector_attribution_results.json}.

\section{Первый цикл: базовый мост и первичный аудит}

\subsection{Настройка}

На первом этапе были взяты параметры
\[
R_3=1,\qquad R_1=1,
\]
с усечением спектра
\[
n_{\max}=80,\qquad m_{\max}=80,
\]
и fit-окном
\[
t\in\{0.03,0.035,0.04,0.045,0.05,0.055,0.06\}.
\]

\subsection{Основные результаты первого цикла}

\begin{center}
\begin{tabular}{lcc}
\toprule
Величина & Численный fit & Теоретическое значение \\
\midrule
$a_0$ & 124.0268873911 & $4\pi^3 = 124.0251067212$ \\
$a_2$ & 123.8993672201 & $4\pi^3 = 124.0251067212$ \\
$a_2/(4\pi)$ & 9.8595983695 & $\pi^2 = 9.8696044011$ \\
\bottomrule
\end{tabular}
\end{center}

Относительные ошибки составили:
\begin{itemize}
    \item для \(a_0\): \(1.44\times 10^{-5}\);
    \item для \(a_2\): \(1.01\times 10^{-3}\);
    \item для \(a_2/(4\pi)\): \(1.01\times 10^{-3}\).
\end{itemize}

\subsection{Интерпретация первого цикла}

Первичный аудит дал три главных результата.
\begin{enumerate}
    \item Ведущий коэффициент теплового следа воспроизводит масштаб \(4\pi^3\) почти идеально.
    \item Нормированный кривизновый вклад воспроизводит масштаб \(\pi^2\) с ошибкой порядка \(10^{-3}\).
    \item Длина фундаментального цикла \(S^1\) даёт естественный систолический proxy
    \[
    \frac{\mathrm{sys}(S^1)}{2}=\pi,
    \]
    что делает линейный \(\pi\)-сектор геометрически правдоподобным, а не чисто нумерологическим.
\end{enumerate}

Именно после этого первого цикла гипотеза ``TOE и UGSM являются двумя тенями одного спектрального источника'' перестала быть просто философской аналогией и получила количественную опору.

\section{Второй цикл: углублённый аудит устойчивости}

Второй цикл был нужен, чтобы проверить более жёсткий вопрос: не является ли найденное совпадение лишь эффектом одной специальной численной настройки.

\subsection{Проверка по окну фита}

Для \(R_3=R_1=1\) были рассмотрены три окна:
\begin{itemize}
    \item narrow: \(0.02\) -- \(0.04\);
    \item baseline: \(0.03\) -- \(0.06\);
    \item wide: \(0.04\) -- \(0.08\).
\end{itemize}

\begin{center}
\begin{tabular}{lccc}
\toprule
Окно & $a_0$ relative error & $a_2$ relative error & $a_2/(4\pi)$ \\
\midrule
narrow & $4.17\times 10^{-6}$ & $4.45\times 10^{-4}$ & 9.8652 \\
baseline & $1.44\times 10^{-5}$ & $1.01\times 10^{-3}$ & 9.8596 \\
wide & $3.41\times 10^{-5}$ & $1.82\times 10^{-3}$ & 9.8517 \\
\bottomrule
\end{tabular}
\end{center}

Картина остаётся устойчивой: \(a_0\) стабилен на уровне \(10^{-5}\), а \(a_2\) --- на уровне \(10^{-3}\). Значит, базовый результат не является артефактом одного специально подобранного окна.

\subsection{Радиусный прогон}

Далее были проверены пары радиусов
\[
(1.0,1.0),\ (1.2,1.0),\ (1.0,1.3),\ (1.5,0.8),\ (0.8,1.5).
\]

Теоретические ожидания для скалярного сектора были взяты в форме
\begin{align}
 a_0^{exp} &= 4\pi^3 R_3^3R_1,\\
 a_2^{exp} &= 4\pi^3 R_3R_1.
\end{align}

\begin{longtable}{p{2.0cm}p{2.0cm}p{2.5cm}p{2.5cm}p{2.4cm}p{2.4cm}}
\toprule
$R_3$ & $R_1$ & $a_0$ rel. err. & $a_2$ rel. err. & $a_2/(4\pi)$ fit & $a_2/(4\pi)$ exp. \\
\midrule
\endhead
1.0 & 1.0 & $1.44\times 10^{-5}$ & $1.01\times 10^{-3}$ & 9.8596 & 9.8696 \\
1.2 & 1.0 & $4.76\times 10^{-6}$ & $4.84\times 10^{-4}$ & 11.8378 & 11.8435 \\
1.0 & 1.3 & $1.44\times 10^{-5}$ & $1.01\times 10^{-3}$ & 12.8175 & 12.8305 \\
1.5 & 0.8 & $1.24\times 10^{-6}$ & $1.97\times 10^{-4}$ & 11.8412 & 11.8435 \\
0.8 & 1.5 & $5.58\times 10^{-5}$ & $2.52\times 10^{-3}$ & 11.8137 & 11.8435 \\
\bottomrule
\end{longtable}

Это подтверждает, что:
\begin{itemize}
    \item объёмный сектор следует масштабу \(R_3^3R_1\) с ошибками порядка \(10^{-6}\) -- \(10^{-5}\);
    \item кривизновый сектор следует масштабу \(R_3R_1\) с ошибками порядка \(10^{-4}\) -- \(10^{-3}\).
\end{itemize}

Следовательно, мост не разрушается при деформации геометрии, а ведёт себя так, как должен вести себя единый геометрический источник.

\subsection{Проверка на семействе спектральных kernels}

Чтобы выйти за пределы одного теплового ядра, были исследованы kernels
\[
e^{-x},\qquad e^{-2x},\qquad (1+x)^{-3},\qquad (1+x)^{-4}.
\]

Для каждого из них вычислялось спектральное действие
\[
S_f(\Lambda)=\mathrm{Tr}\,f(D^2/\Lambda^2)
\]
при \(\Lambda=6,8,10,12\), а затем анализировалась объёмно-нормированная комбинация
\[
\frac{S_f(\Lambda)}{\Lambda^4R_3^3R_1}.
\]

\begin{center}
\begin{tabular}{lcc}
\toprule
Kernel & Mean normalized value & Relative spread on $\Lambda=6\ldots12$ \\
\midrule
$e^{-x}$ & 0.79736 & 2.09\% \\
$e^{-2x}$ & 0.20239 & 4.19\% \\
$(1+x)^{-3}$ & 0.39343 & 3.83\% \\
$(1+x)^{-4}$ & 0.13494 & 4.30\% \\
\bottomrule
\end{tabular}
\end{center}

Этот тест не требует совпадения абсолютных нормировок для разных kernels; вместо этого он показывает, что разные функциональные окна продолжают видеть один и тот же объёмно-геометрический каркас. Это важный аргумент в пользу того, что TOE и UGSM могут быть не двумя разными законами, а двумя чтениями одного и того же спектра.


\section{Третий цикл: Dirac-type перенос моста}

После двух скалярных циклов нужно было проверить, сохраняется ли найденная структура при переходе к более физически содержательному оператору. Для этого мост был перенесён в Dirac-type сектор, где рассматривался квадрат минимального спинорного оператора на \(S^3\times S^1\) с периодической spin-структурой вдоль окружности.

\subsection{Постановка}

В третьем цикле использовался спектр вида
\[
\lambda_{n,m}^{(D^2)}=\frac{(n+3/2)^2}{R_3^2}+\frac{m^2}{R_1^2},
\qquad m\in\mathbb{Z},
\]
с кратностями
\[
g_n^{(D^2)}=2(n+1)(n+2).
\]

Как и раньше, анализировался нормированный тепловой след
\[
\widetilde K_D(t)=(4\pi t)^2\,\mathrm{Tr}\,e^{-tD^2},
\]
который аппроксимировался в виде
\[
\widetilde K_D(t)\approx a_{0,D}+a_{2,D}t+a_{4,D}t^2.
\]

Для минимального Dirac-type сектора теоретические ожидания принимают форму
\begin{align}
a_{0,D}^{exp} &= 8\pi^3 R_3^3R_1,\\
a_{2,D}^{exp} &= -4\pi^3 R_3R_1,
\end{align}
так что нормированная субведущая комбинация должна воспроизводить
\[
\frac{a_{2,D}}{4\pi}=-\pi^2 R_3R_1.
\]

\subsection{Основные результаты третьего цикла}

Для базовой точки \(R_3=R_1=1\), усечения \(n_{\max}=m_{\max}=100\) и baseline-окна \(t=0.03\) -- \(0.06\) получаются значения
\begin{center}
\begin{tabular}{lcc}
\toprule
Величина & Численный fit & Теоретическое значение \\
\midrule
$a_{0,D}$ & 248.0502134424 & $8\pi^3 = 248.0502134424$ \\
$a_{2,D}$ & -124.0251067212 & $-4\pi^3 = -124.0251067212$ \\
$a_{2,D}/(4\pi)$ & -9.8696044011 & $-\pi^2 = -9.8696044011$ \\
\bottomrule
\end{tabular}
\end{center}

Относительные ошибки оказываются исключительно малы:
\begin{itemize}
    \item для \(a_{0,D}\): \(1.72\times 10^{-15}\);
    \item для \(a_{2,D}\): \(2.27\times 10^{-13}\).
\end{itemize}

Проверка по окнам фита показывает ту же картину. Для narrow, baseline и wide окон ошибки по \(a_{0,D}\) остаются на уровне \(10^{-15}\), а ошибки по \(a_{2,D}\) --- на уровне \(10^{-13}\). Это существенно сильнее, чем в скалярном секторе, и означает, что Dirac-type мост не просто существует, а численно фиксируется почти с машинной точностью.

Радиусный прогон также подтверждает ожидаемое масштабирование. При парах \((R_3,R_1)=(1.2,1.0)\), \((1.0,1.3)\), \((1.5,0.8)\), \((0.8,1.5)\) коэффициент \(a_{0,D}\) следует закону \(R_3^3R_1\), а коэффициент \(a_{2,D}\) --- закону \(R_3R_1\), причём типичные относительные ошибки лежат в диапазоне \(10^{-16}\) -- \(10^{-13}\).

Наконец, проверка на семействе kernels показывает, что разные функциональные окна продолжают читать один и тот же объёмно-нормированный Dirac-type каркас. Для \(e^{-x}\), \(e^{-2x}\), \((1+x)^{-3}\) и \((1+x)^{-4}\) относительный разброс нормированной величины по \(\Lambda=6,8,10,12\) остаётся на уровне примерно \(0.8\%\) -- \(2.1\%\).

\subsection{Вывод третьего цикла}

Третий цикл показывает, что мост TOE--UGSM не привязан к одному скалярному toy-sector. При переходе к минимальному спинорному Dirac-type оператору сохраняются те же геометрические масштабы, но уже в более физически содержательной постановке: ведущий коэффициент удваивается до \(8\pi^3\), кривизновый вклад принимает знак, ожидаемый для Dirac-type сектора, а нормированная субведущая часть точно воспроизводит \(-\pi^2\). Именно после этого шага становится естественно переходить к четвёртому циклу и спрашивать уже не о том, существует ли сам спинорный мост, а о том, как в нём устроен глобальный holonomy/\(\pi\)-сектор.


\section{Четвёртый цикл: spin-holonomy анализ}

После Dirac-type цикла оставался главный незакрытый вопрос: действительно ли линейный \(\pi\)-сектор связан с глобальной структурой цикла \(S^1\), а не просто с внешней proxy-интерпретацией. Для этого был введён twist-параметр \(\alpha\in[0,1/2]\), который одновременно кодирует spin-структуру и Wilson/holonomy-сдвиг вдоль окружности.

\subsection{Постановка}

В spin-holonomy цикле спектр на окружности заменяется на
\[
\lambda_{S^1}(m,\alpha)=\frac{(m+\alpha)^2}{R_1^2},
\qquad m\in\mathbb{Z},
\]
где:
\begin{itemize}
    \item \(\alpha=0\) соответствует периодической spin-структуре;
    \item \(\alpha=1/2\) соответствует анти-периодической spin-структуре;
    \item промежуточные значения \(\alpha\) интерпретируются как twisted holonomy/Wilson-line деформации.
\end{itemize}

Соответствующая Wilson-фаза равна
\[
\theta=2\pi\alpha.
\]

Именно здесь появляется ключевой механизм: при \(\alpha=1/2\) автоматически получается
\[
\theta=\pi.
\]

Это уже не внешняя подстановка, а прямой геометрический вывод линейного \(\pi\)-сектора как половинной holonomy-фазы вокруг фундаментального цикла.

\subsection{Что проверялось}

Были проверены значения
\[
\alpha\in\{0,0.125,0.25,0.375,0.5\}.
\]

Для каждого \(\alpha\) вычислялись:
\begin{itemize}
    \item коэффициенты \(a_{0,D}\) и \(a_{2,D}\) из Dirac-type теплового следа;
    \item Wilson-фаза \(\theta=2\pi\alpha\);
    \item минимальный зазор кругового сектора
    \[
    \Delta_{S^1}(\alpha)=\frac{\alpha}{R_1},
    \]
    который измеряет чувствительность глобального цикла к spin/holonomy-сдвигу.
\end{itemize}

\subsection{Главный численный результат}

Результат оказался исключительно сильным: локальные коэффициенты \(a_{0,D}\) и \(a_{2,D}\) остаются неизменными при изменении \(\alpha\) вплоть до машинной точности, тогда как глобальная фаза и круговой gap меняются линейно с \(\alpha\).

\begin{center}
\begin{tabular}{lccccc}
\toprule
$\alpha$ & $\theta/\pi$ & $a_{0,D}$ rel. err. & $a_{2,D}$ rel. err. & $a_{2,D}/(4\pi)$ & gap on $S^1$ \\
\midrule
0.000 & 0.00 & $5.73\times 10^{-16}$ & $1.20\times 10^{-13}$ & -9.8696 & 0.000 \\
0.125 & 0.25 & $2.06\times 10^{-15}$ & $1.43\times 10^{-13}$ & -9.8696 & 0.125 \\
0.250 & 0.50 & $5.73\times 10^{-16}$ & $1.03\times 10^{-13}$ & -9.8696 & 0.250 \\
0.375 & 0.75 & $2.06\times 10^{-15}$ & $2.46\times 10^{-13}$ & -9.8696 & 0.375 \\
0.500 & 1.00 & $1.60\times 10^{-15}$ & $1.75\times 10^{-13}$ & -9.8696 & 0.500 \\
\bottomrule
\end{tabular}
\end{center}

Сводные метрики инвариантности дают:
\begin{itemize}
    \item максимальная относительная ошибка по \(a_{0,D}\): \(2.06\times 10^{-15}\);
    \item максимальная относительная ошибка по \(a_{2,D}\): \(2.46\times 10^{-13}\);
    \item полный разброс по \(a_{0,D}\): около \(10^{-12}\);
    \item полный разброс по \(a_{2,D}\): около \(5\times 10^{-11}\).
\end{itemize}

\subsection{Интерпретация четвёртого цикла}

Этот цикл впервые отделяет два слоя физико-геометрической структуры:
\begin{enumerate}
    \item \textbf{Локальный слой.} Коэффициенты Seeley--DeWitt, а значит объёмный и кривизновый сектор, не зависят от holonomy-сдвига. Это именно то, что и должно происходить для локальных инвариантов.
    \item \textbf{Глобальный слой.} Wilson-фаза и минимальный круговой gap чувствительны к \(\alpha\), то есть к spin-структуре и глобальной обмотке вдоль \(S^1\).
\end{enumerate}

Отсюда следует качественно новый вывод: линейный \(\pi\)-сектор действительно может быть связан не с локальной частью тепловой асимптотики, а с глобальным циклическим holonomy-сектором. Более того, при анти-периодической spin-структуре это значение возникает строго:
\[
\theta_{AP}=2\pi\cdot \frac{1}{2}=\pi.
\]

Это существенно сильнее прежнего языка про ``proxy'': теперь \(\pi\) появляется как прямой геометрический след half-twist/half-holonomy на фундаментальном цикле.

\subsection{Что добавляет четвёртый цикл к гипотезе единства}

Четвёртый цикл усиливает статью сразу в двух отношениях.
\begin{itemize}
    \item Он показывает, что локальные коэффициенты и глобальный \(\pi\)-сектор можно строго развести по разным уровням одного и того же операторного описания.
    \item Он даёт конкретный механизм, как \(\pi\)-сектор может возникать из spin/holonomy-структуры на \(S^1\), а значит перестаёт быть просто эвристическим геометрическим комментарием.
\end{itemize}

Именно это делает гипотезу общего спектрального источника не только устойчивой и физически содержательной, но и значительно более структурно объяснённой.




\section{Пятый цикл: gauge-holonomy расширение}

Следующий шаг после spin-holonomy анализа состоял в минимальном переходе от чистого twist-параметра \(\alpha\) к более полной Wilson-loop/gauge-holonomy деформации. Для этого был введён внутренний doublet-сектор с зарядами
\[
q=\pm 1,
\]
и дополнительный абелев holonomy-параметр \(\beta\). Тогда круговой спектр принимает вид
\[
\lambda_{S^1}(m;\alpha,\beta,q)=\frac{(m+\alpha+q\beta)^2}{R_1^2}.
\]

В этом цикле фиксировалась анти-периодическая spin-структура \(\alpha=1/2\), а затем исследовалась gauge-holonomy деформация \(\beta\in\{0,0.125,0.25,0.375,0.5\}\).

\subsection{Что именно проверялось}

Проверялись две вещи:
\begin{enumerate}
    \item сохраняются ли локальные Dirac-type коэффициенты при включении gauge-holonomy doublet-сектора;
    \item как именно перестраиваются глобальные Wilson-фазы и эффективный круговой gap при изменении \(\beta\).
\end{enumerate}

Введённые фазы для двух зарядовых ветвей равны
\[
\theta_\pm = 2\pi(\alpha\pm\beta).
\]

Ожидаемые локальные коэффициенты просто удваиваются относительно одного Dirac-doublet сектора:
\begin{align}
 a_{0,g}^{exp} &= 16\pi^3 R_3^3R_1,\\
 a_{2,g}^{exp} &= -8\pi^3 R_3R_1,
\end{align}
а значит
\[
\frac{a_{2,g}^{exp}}{4\pi}=-2\pi^2.
\]

\subsection{Численный результат}

Для всех исследованных \(\beta\) локальные коэффициенты остаются неизменными до машинной точности, тогда как глобальные фазы перестраиваются линейно.

\begin{center}
\begin{tabular}{lcccccc}
\toprule
$\beta$ & $\theta_+/\pi$ & $\theta_-/\pi$ & $a_{0,g}$ rel. err. & $a_{2,g}$ rel. err. & $a_{2,g}/(4\pi)$ & eff. gap \\
\midrule
0.000 & 1.00 & 1.00 & $9.17\times 10^{-16}$ & $1.15\times 10^{-13}$ & -19.7392 & 0.500 \\
0.125 & 1.25 & 0.75 & $1.60\times 10^{-15}$ & $1.81\times 10^{-13}$ & -19.7392 & 0.375 \\
0.250 & 1.50 & 0.50 & $1.15\times 10^{-16}$ & $5.34\times 10^{-14}$ & -19.7392 & 0.250 \\
0.375 & 1.75 & 0.25 & $1.03\times 10^{-15}$ & $1.50\times 10^{-13}$ & -19.7392 & 0.125 \\
0.500 & 2.00 & 0.00 & $1.49\times 10^{-15}$ & $8.43\times 10^{-14}$ & -19.7392 & 0.000 \\
\bottomrule
\end{tabular}
\end{center}

\begin{figure}[ht]
\centering
\includegraphics[width=0.78\textwidth]{fig_gauge_holonomy_phases.png}
\caption{Gauge-holonomy цикл: фазовые ветви \(\theta_\pm\) и эффективный круговой gap как функции \(\beta\).}
\end{figure}

\subsection{Вывод пятого цикла}

Gauge-holonomy расширение не ломает локальное спектральное ядро, но делает глобальный фазовый сектор богаче: одна фаза \(\theta\) превращается в семейство ветвей \(\theta_\pm\). Линейный \(\pi\)-сектор при этом не исчезает, а становится специальной точкой внутри более общей holonomy-структуры.

\section{Шестой цикл: enriched Dirac-type оператор}

После gauge-holonomy цикла следующий естественный переход состоял в построении более богатого оператора с несколькими внутренними каналами. В шестом цикле был рассмотрен минимальный enriched Dirac-type оператор, в котором имеются четыре внутренние ветви:
\[
(q,\mu)\in\{(+1,0),(-1,0),(+1,0.6),(-1,0.6)\}.
\]

Это минимальный внутренний charge/mass multiplet: два лёгких канала и два тяжёлых канала, каждый со своей holonomy-фазой. Соответствующий спектр брался в виде
\[
\lambda_{n,m}^{(rich)}=\frac{(n+3/2)^2}{R_3^2}+\frac{(m+\alpha+q\beta)^2}{R_1^2}+\mu^2.
\]

В базовом прогоне фиксировались
\[
R_3=R_1=1,\qquad \alpha=\frac{1}{2},\qquad \beta=0.25.
\]

Тогда две фазовые ветви равны
\[
\theta_+=\frac{3\pi}{2},
\qquad
\theta_-=\frac{\pi}{2}.
\]

\subsection{Численный результат}

Расчёт дал следующие значения:
\begin{center}
\begin{tabular}{lcc}
\toprule
Величина & Численный результат & Эффективная модель / комментарий \\
\midrule
$a_{0,rich}$ & 992.1992 & 992.2009 \\
$a_{2,rich}$ & -674.5806 & -674.6966 \\
$a_{2,rich}/(4\pi)$ & -53.6814 & subleading enriched sector \\
Effective gap & 0.25 & controlled by $\alpha\pm\beta$ branches \\
\bottomrule
\end{tabular}
\end{center}

Ошибки остаются малыми: для \(a_{0,rich}\) --- \(1.65\times 10^{-6}\), а для сравнения \(a_{2,rich}\) с эффективной enriched-моделью --- \(1.72\times 10^{-4}\).

Внутренний баланс секторов при \(t=0.04\) почти симметричен:
\begin{center}
\begin{tabular}{lc}
\toprule
Сектор & Доля в тепловом следе \\
\midrule
$q_+^{light}$ & 0.2518 \\
$q_-^{light}$ & 0.2518 \\
$q_+^{heavy}$ & 0.2482 \\
$q_-^{heavy}$ & 0.2482 \\
\bottomrule
\end{tabular}
\end{center}

\subsection{Вывод шестого цикла}

Мост выдерживает переход от чисто геометрического ядра к richer оператору, в котором subleading-сектор уже несёт информацию о внутренней mass-структуре. Ведущий объёмный сектор остаётся под контролем, а глобальные фазы \(\theta_\pm\) продолжают задавать holonomy-ветви внутри многоканального оператора.

\section{Седьмой цикл: sector-map анализ}

После построения enriched оператора следующий естественный переход состоял в отделении двух эффектов, которые в шестом цикле были ещё смешаны: фазового holonomy-сектора и внутренней mass-driven subleading-структуры. Для этого был проведён отдельный sector-map анализ по параметрам
\[
\beta\in\{0,0.125,0.25,0.375,0.5\},
\qquad
\mu_{heavy}\in\{0,0.3,0.6,0.9\}.
\]

\subsection{Главный результат карты параметров}

Максимальные ошибки по всей карте остаются малыми:
\begin{itemize}
    \item максимальная ошибка по \(a_0\): \(1.02\times 10^{-5}\);
    \item максимальная ошибка по \(a_2\) относительно эффективной enriched-модели: \(7.99\times 10^{-4}\).
\end{itemize}

Средний subleading коэффициент по срезам \(\mu_{heavy}\) ведёт себя так:
\begin{center}
\begin{tabular}{lcc}
\toprule
$\mu_{heavy}$ & Mean $a_2$ fit & Mean $a_2/(4\pi)$ \\
\midrule
0.0 & -496.1004 & -39.4784 \\
0.3 & -540.7432 & -43.0310 \\
0.6 & -674.5806 & -53.6814 \\
0.9 & -897.2242 & -71.3988 \\
\bottomrule
\end{tabular}
\end{center}

Напротив, усреднение по срезам \(\beta\) показывает почти полную инвариантность среднего \(a_2\):
\begin{center}
\begin{tabular}{lccc}
\toprule
$\beta$ & $\theta_+/\pi$ & $\theta_-/\pi$ & Mean $a_2$ fit \\
\midrule
0.000 & 1.00 & 1.00 & -652.1621 \\
0.125 & 1.25 & 0.75 & -652.1621 \\
0.250 & 1.50 & 0.50 & -652.1621 \\
0.375 & 1.75 & 0.25 & -652.1621 \\
0.500 & 2.00 & 0.00 & -652.1621 \\
\bottomrule
\end{tabular}
\end{center}

\begin{figure}[ht]
\centering
\includegraphics[width=0.72\textwidth]{fig_enriched_heatmap.png}
\caption{Heatmap для \(a_2/(4\pi)\) в enriched операторе на карте \((\mu_{heavy},\beta)\). Видно, что карта чувствительна прежде всего к \(\mu_{heavy}\), а не к \(\beta\).}
\end{figure}

\subsection{Вывод седьмого цикла}

В richer операторе \(\beta\) перестраивает фазовую конфигурацию \(\theta_\pm\), но почти не влияет на средний subleading коэффициент, тогда как внутренний mass-параметр \(\mu_{heavy}\) существенно управляет subleading спектральной структурой. Это уже почти факторизация ролей внутри оператора.

\section{Восьмой цикл: sector-attribution анализ}

После sector-map цикла стало возможно сделать следующий качественный переход: не просто видеть, что разные параметры влияют на разные части спектральной структуры, а сформулировать явный словарь соответствий между компонентами richer оператора и теми смысловыми блоками, которые потенциально могут быть связаны с TOE и UGSM.

\subsection{Sensitivity-метрики}

На основе карты \((\mu,\beta)\) были вычислены средние чувствительности для subleading коэффициента \(a_2\):
\begin{itemize}
    \item
    \[
    \left\langle \left|\frac{d a_2}{d\mu}\right| \right\rangle \approx 445.6931;
    \]
    \item
    \[
    \left\langle \left|\frac{d a_2}{d\beta}\right| \right\rangle \approx 2.53\times 10^{-10}.
    \]
\end{itemize}

Их отношение даёт separation ratio
\[
\mathcal{R}_{sep}\approx 1.76\times 10^{12}.
\]

\begin{figure}[ht]
\centering
\includegraphics[width=0.66\textwidth]{fig_sector_attribution.png}
\caption{Разделение чувствительностей: subleading коэффициент \(a_2\) на много порядков сильнее зависит от \(\mu\), чем от \(\beta\).}
\end{figure}

\subsection{Рабочий словарь соответствий}

В минимальной рабочей форме operator dictionary выглядит так:
\begin{center}
\begin{tabular}{p{3.6cm}p{3.8cm}p{3.3cm}p{3.3cm}}
\toprule
Компонента оператора & Роль в enriched-модели & TOE-чтение & UGSM-чтение \\
\midrule
Ведущий член $a_0$ & геометрический vacuum backbone & background spectral-correlation sector & vacuum geometric normalization \\
Subleading член $a_2$ & локальная curvature + internal load & local response block & EFT curvature/mass-sensitive sector \\
Фазы $\theta_\pm$ & глобальный holonomy/Wilson сектор & modular/global phase branch & compact-cycle phase and $\pi$-sector \\
Heavy-channel $\mu$ & внутренняя subleading нагрузка & hidden/internal spectral substructure & internal multiplet deformation of observables \\
\bottomrule
\end{tabular}
\end{center}

\subsection{Вывод восьмого цикла}

Программа переходит из режима ``численные совпадения + интерпретации'' в режим ``операторной атрибуции ролей''. Это ещё не глобальная теорема об эквивалентности, но уже каркас для будущей sector-reduction theorem: есть устойчивый геометрический backbone, отдельно выделенный global-phase sector, отдельно выделенный internal-load sector и количественная метрика, что эти роли действительно различимы.

\section{Визуальные суммарные проверки}

Для полноты полезно собрать визуально три ключевых наблюдения статьи.

\begin{figure}[ht]
\centering
\includegraphics[width=0.66\textwidth]{fig_scalar_window_stability.png}
\caption{Скалярный цикл: устойчивость относительных ошибок \(a_0\) и \(a_2\) при изменении окна фита.}
\end{figure}

Вместе рисунки показывают общую иерархию: геометрический backbone стабилен, holonomy-параметры управляют глобальными фазами, а внутренние mass-channel'ы управляют subleading-структурой.

\section{Синтез восьми циклов}

После объединения результатов восьми циклов картина выглядит значительно сильнее, чем после первого моста.

\subsection{Что показали циклы 1--4}
\begin{itemize}
    \item минимальный мост между TOE и UGSM существует и воспроизводит масштабы \(4\pi^3\), \(\pi^2\) и \(\pi\);
    \item картина устойчива к смене окна фита и деформации радиусов;
    \item мост переносится в Dirac-type сектор;
    \item линейный \(\pi\)-сектор в минимальном spin-holonomy секторе возникает как строгое значение holonomy-фазы \(\theta=\pi\).
\end{itemize}

\subsection{Что показали циклы 5--8}
\begin{itemize}
    \item gauge-holonomy расширение не ломает локальные коэффициенты и превращает одну фазу в семейство ветвей \(\theta_\pm\);
    \item richer внутренний operator/multiplet сохраняет ведущий спектральный каркас и добавляет внутреннюю subleading-нагрузку;
    \item sector-map анализ отделяет mass-driven subleading сектор от holonomy-driven фазового сектора;
    \item sector-attribution цикл вводит первую явную схему соответствий между операторными компонентами и возможными блоками TOE/UGSM.
\end{itemize}

\subsection{Общий вывод из объединения}

Совместно восемь циклов показывают следующее: гипотеза единого спектрального источника демонстрирует структурную устойчивость, переносимость в более физически содержательные операторные сектора, естественный механизм возникновения \(\pi\)-сектора из spin/holonomy-геометрии, устойчивость при gauge-holonomy расширении и первую рабочую схему operator-attribution. Для исследовательского уровня это уже не просто серия удачных совпадений, а устойчивая и постепенно структурируемая программа секторной редукции.

\section{Ограничения исследования}

Несмотря на силу результата, статья не утверждает полного доказательства единства TOE и UGSM. Сохраняются следующие ограничения:
\begin{itemize}
    \item уже проверен минимальный Dirac-type сектор, но ещё не исчерпан полный класс физически релевантных Dirac-операторов и spin-структур;
    \item уже построена минимальная enriched sector-attribution схема, но ещё не доказано, что этот словарь единственен и что он продолжается до полного оператора с действительно богатой внутренней калибровочной структурой и динамикой между каналами;
    \item не доказано, что полный корреляционный оператор TOE глобально совпадает с операторным источником UGSM, а не только редуцируется к нему в компактном секторе;
    \item линейный \(\pi\)-сектор уже имеет строгий вывод в минимальном spin-holonomy секторе, но ещё не выведена его полная феноменологическая роль в более богатой gauge-внутренней структуре;
    \item мост пока подтверждён как исследовательская программа высокого доверия, но не как завершённая редукция одной теории к другой.
\end{itemize}

\section{Обсуждение}

С методологической точки зрения главный результат статьи состоит в том, что обе теории помещены в один общий режим проверки, а сам проверяемый объект постепенно эволюционирует: от скалярного теплового следа к spin-holonomy сектору, затем к gauge-holonomy, richer operator и, наконец, к явной атрибуции ролей отдельных компонент. TOE в таком чтении поставляет фундаментальный корреляционный источник, UGSM --- дисциплинированное спектральное чтение компактного сектора, а многослойный аудит связывает эти уровни через конкретный геометрический и операторный след.

\section{Заключение}

На основании объединённого исследования можно сформулировать итоговый вывод статьи.

\begin{quote}
На текущем этапе \texttt{TOE.pdf} и \texttt{ugsm\_dynamics\_audit-3.pdf} правдоподобно интерпретируются как две разные тени одной более общей спектральной сущности. Этот вывод опирается на восемь последовательных циклов проверки: от базового скалярного моста к Dirac-type сектору, затем к spin- и gauge-holonomy расширениям, далее к richer внутреннему оператору, карте секторов и, наконец, к явной operator-attribution схеме.
\end{quote}

Самая аккуратная формулировка такова: единство двух теорий ещё не доказано окончательно, но гипотеза общего спектрального источника после проведённого исследования стала существенно более конкретной, более устойчивой и заметно более структурированной, чем на исходном уровне одной численной аналогии.

\section{Дальнейшая программа}

После восьми последовательных циклов проверки речь уже идёт не о ближайшем ``следующем шаге'' в прежнем смысле, а о следующем исследовательском слое. Теперь наиболее естественно двигаться в трёх связанных направлениях:
\begin{itemize}
    \item перейти от минимального абелевого holonomy-сектора к более богатым Wilson-loop и gauge-структурам, где внутренние степени свободы образуют уже не просто doublet, а действительно нетривиальную алгебру;
    \item построить более полный Dirac-type оператор с внутренней динамикой, в котором каналы TOE- и UGSM-типа связываются не только через статические mass-loadings, но и через операторные взаимодействия;
    \item перевести нынешнюю sector-map/sector-attribution схему в режим строгих утверждений и доказуемых критериев sector reduction, то есть точно описать, какие блоки richer оператора допускают устойчивое отождествление с наблюдаемыми структурами TOE и UGSM.
\end{itemize}

Именно этот слой должен решить, может ли восьмикратно проверенная рабочая гипотеза быть поднята от сильной исследовательской интерпретации до существенно более строгой программы доказательства единства.

\end{document}